%
% mypreface.tex
%

% put the following in your doc.tex
%\usepackage{./code/mips}
%\input{./code/mips.sty}

% lauracode.txt BEGIN
% ******************************************************************
% *** SPACE CONTROL ************************************************

\def\vf{\vfill}
\def\vs{\vspace{1pc}}
\def\hs{\hspace{1in}}

\def\blank{$\,$ \underline{\hspace{1in}} $\,$}
\def\Blank{$\,$ \underline{\hspace{1.5in}} $\,$}

\def\pg{\clearpage}

% ******************************************************************
% *** FORMATTING ***************************************************

\def\und{\underline}

% ******************************************************************
% *** ENGLISH ABBREVIATIONS ****************************************

\def\ie{\mbox{{\it i.e.} }}
\def\eg{\mbox{{\it e.g.} }}
\def\etc{\mbox{{\it et cetera} }}

% ******************************************************************
% *** MATH SHORTCUTS ***********************************************

\def\ds{\displaystyle}

\def\ddx{\ds \frac{d}{dx}}
\def\ddt{\ds \frac{d}{dx}}
\newcommand{\dd}[2]{\ds \frac{d{#1}}{d{#2}}}

\def\pitch{\mathbin{\frown \! \! \! \! \mid \, \, \,}}  % replace by \pitchfork if amssymb is working

\def\qed{\rule[-.23ex]{1.6ex}{1.6ex}}

% ******************************************************************
% *** TABLE SHORTHAND **********************************************

\newcommand{\gotable}[3]
      {\begin{table}[h] 
       \begin{center}
       \addtolength{\tabcolsep}{#1mm}
       \renewcommand{\arraystretch}{#2}
       \begin{tabular}{#3}}

\newcommand{\stoptable}
      {\end{tabular}
       \end{center}
       \end{table}}
       
% ******************************************************************
% *** MINIPAGE SHORTHAND *******************************************

\newcommand{\twofig}[4]
   {\begin{figure}[h]
      \begin{minipage}[t]{3in}
        \begin{center}
        {\Large \bf (#1) $\!\!\!\!\!\!\!\!\!$}
        \epsfig{file=#2,height=2in,width=2in} 
        \end{center}
      \end{minipage}
      \hfill
      \begin{minipage}[t]{3in}
        \begin{center}
        {\Large \bf (#3) $\!\!\!\!\!\!\!\!\!$}
        \epsfig{file=#4,height=2in,width=2in}  
        \end{center}
      \end{minipage}
    \end{figure}}

\newcommand{\threefig}[6]
   {\begin{figure}[h]
      \begin{minipage}[t]{2in}
        \begin{center}
        {\large \bf (#1) $\!\!\!\!\!\!\!\!\!$}
        \epsfig{file=#2,height=1.6in,width=1.6in}
        \end{center}
      \end{minipage}
      \hfill
      \begin{minipage}[t]{2in}
        \begin{center}
        {\large \bf (#3) $\!\!\!\!\!\!\!\!\!$}
        \epsfig{file=#4,height=1.6in,width=1.6in}
        \end{center}
      \end{minipage}
      \hfill
      \begin{minipage}[t]{2in}
        \begin{center}
        {\large \bf (#5) $\!\!\!\!\!\!\!\!\!$}
        \epsfig{file=#6,height=1.6in,width=1.6in}
        \end{center}
      \end{minipage}
    \end{figure}}
         
\newcommand{\twoup}[2]
   {\begin{minipage}{3in}
      \begin{center}
      #1 
      \end{center}
    \end{minipage}
    \hfill
    \begin{minipage}{3in}
      \begin{center}
      #2 
      \end{center}
    \end{minipage}}         

% ******************************************************************
% *** TESTPOINTS SHORTHAND *****************************************

\def\bpz{\begin{problem}{0}}
\newcommand{\bp}[1]{\begin{problem}{#1}}

\def\npz{\newpart{0}}
\newcommand{\np}[1]{\newpart{#1}}

\def\ep{\end{problem}}

%\newcommand{\probpart}...make later
\newcommand{\prob}[1]{\setcounter{problemnum}{#1} \Large \bf \theproblemnum .}

% ***********************************************************************
% *** COMMUTATIVE DIAGRAMS **********************************************

\def\normalbaselines{\baselineskip20pt \lineskip3pt \lineskiplimit3pt}

% for extending arrows
\def\ext{\rule[.5ex]{5mm}{.1mm}}   
\def\Ext{\rule[.5ex]{10mm}{.1mm}}  
\def\back{\!\!\!}
\def\Back{\back\back}

% right maps
\def\mapR{\smash{\mathop{\longrightarrow}}}
\def\mapRR{\smash{\mathop{\ext\back\longrightarrow}}}
\def\mapRRR{\smash{\mathop{\Ext\back\longrightarrow}}}

\def\MapR#1#2{\smash{\mathop{\longrightarrow}\limits^{#1}_{#2}}}
\def\MapRR#1#2{\smash{\mathop{\ext\back\longrightarrow}\limits^{#1}_{#2}}}
\def\MapRRR#1#2{\smash{\mathop{\Ext\back\longrightarrow}\limits^{#1}_{#2}}}

% right injective maps
\def\inprefix{\mbox{\raisebox{1mm}{\tiny $\subset$}}\!\!}

\def\mapinR{\smash{\mathop{\inprefix\!\!\longrightarrow}}}
\def\mapinRR{\smash{\mathop{\inprefix\ext\back\longrightarrow}}}
\def\mapinRRR{\smash{\mathop{\inprefix\Ext\back\longrightarrow}}}

\def\MapinR#1#2{\smash{\mathop{\inprefix\!\!\longrightarrow}\limits^{#1}_{#2}}}
\def\MapinRR#1#2{\smash{\mathop{\inprefix\ext\back\longrightarrow}\limits^{#1}_{#2}}}
\def\MapinRRR#1#2{\smash{\mathop{\inprefix\Ext\back\longrightarrow}\limits^{#1}_{#2}}}

% right surjective maps
\def\onsuffix{\longrightarrow\back\back\!\rightarrow}

\def\maponR{\smash{\mathop{\onsuffix}}}
\def\maponRR{\smash{\mathop{\ext\back\onsuffix}}}
\def\maponRRR{\smash{\mathop{\Ext\back\onsuffix}}}

\def\MaponR#1#2{\smash{\mathop{\onsuffix}\limits^{#1}_{#2}}}
\def\MaponRR#1#2{\smash{\mathop{\ext\back\onsuffix}\limits^{#1}_{#2}}}
\def\MaponRRR#1#2{\smash{\mathop{\Ext\back\onsuffix}\limits^{#1}_{#2}}}

% right equals maps
\def\smalleq{\rlap{\rule[.3ex]{5mm}{.1mm}}\rule[.7ex]{5mm}{.1mm}}
\def\medeq{\rlap{\rule[.3ex]{10mm}{.1mm}}\rule[.7ex]{10mm}{.1mm}}
\def\largeeq{\rlap{\rule[.3ex]{15mm}{.1mm}}\rule[.7ex]{15mm}{.1mm}}

\def\mapeqR{\smash{\mathop{\smalleq}}}
\def\mapeqRR{\smash{\mathop{\medeq}}}
\def\mapeqRRR{\smash{\mathop{\largeeq}}}

\def\MapeqR#1#2{\smash{\mathop{\smalleq}\limits^{#1}_{#2}}}
\def\MapeqRR#1#2{\smash{\mathop{\medeq}\limits^{#1}_{#2}}}
\def\MapeqRRR#1#2{\smash{\mathop{\largeeq}\limits^{#1}_{#2}}}

% left maps
\def\mapL{\smash{\mathop{\longleftarrow}}}
\def\mapLL{\smash{\mathop{\longleftarrow\back\ext}}}
\def\mapLLL{\smash{\mathop{\longleftarrow\back\Ext}}}

\def\MapL#1#2{\smash{\mathop{\longleftarrow}\limits^{#1}_{#2}}}
\def\MapLL#1#2{\smash{\mathop{\longleftarrow\back\ext}\limits^{#1}_{#2}}}
\def\MapLLL#1#2{\smash{\mathop{\longleftarrow\back\Ext}\limits^{#1}_{#2}}}

% left injective maps
\def\insuffix{\!\!\mbox{\raisebox{1mm}{\tiny $\supset$}}}

\def\mapinL{\smash{\mathop{\longleftarrow\back\insuffix}}}
\def\mapinLL{\smash{\mathop{\longleftarrow\back\ext\insuffix}}}
\def\mapinLLL{\smash{\mathop{\longleftarrow\back\Ext\insuffix}}}

\def\MapinL#1#2{\smash{\mathop{\longleftarrow\back\insuffix}\limits^{#1}_{#2}}}
\def\MapinLL#1#2{\smash{\mathop{\longleftarrow\back\ext\insuffix}\limits^{#1}_{#2}}}
\def\MapinLLL#1#2{\smash{\mathop{\longleftarrow\back\Ext\insuffix}\limits^{#1}_{#2}}}

% left surjective maps
\def\onprefix{\leftarrow\back\back\!\longleftarrow}

\def\maponL{\smash{\mathop{\onprefix}}}
\def\maponLL{\smash{\mathop{\onprefix\back\ext}}}
\def\maponLLL{\smash{\mathop{\onprefix\back\Ext}}}

\def\MaponL#1#2{\smash{\mathop{\onprefix}\limits^{#1}_{#2}}}
\def\MaponLL#1#2{\smash{\mathop{\onprefix\back\ext}\limits^{#1}_{#2}}}
\def\MaponLLL#1#2{\smash{\mathop{\onprefix\back\Ext}\limits^{#1}_{#2}}}

% left equals maps                                % (same as right)
\def\mapeqL{\smash{\mathop{\smalleq}}}
\def\mapeqLL{\smash{\mathop{\medeq}}}
\def\mapeqLLL{\smash{\mathop{\largeeq}}}

\def\MapeqL#1#2{\smash{\mathop{\smalleq}\limits^{#1}_{#2}}}
\def\MapeqLL#1#2{\smash{\mathop{\medeq}\limits^{#1}_{#2}}}
\def\MapeqLLL#1#2{\smash{\mathop{\largeeq}\limits^{#1}_{#2}}}

% down maps
\def\mapD{\Big\downarrow}
\def\MapD#1#2{\Big\downarrow\llap{$\vcenter{\hbox{$\Back\!\scriptstyle#1$}}$}
                            \rlap{$\vcenter{\hbox{$\scriptstyle#2$}}$}}

% down injective maps
\def\inabove{\rlap{\rule{-.15mm}{0mm}\mbox{\raisebox{3.5mm}{\tiny $\cap$}}}}
\def\mapinD{\inabove\Big\downarrow}
\def\MapinD#1#2{\inabove\Big\downarrow
             \llap{$\vcenter{\hbox{$\Back\!\scriptstyle#1$}}$}
             \rlap{$\vcenter{\hbox{$\scriptstyle#2$}}$}}

% down surjective maps
\def\onbelow{\rlap{\rule{.33mm}{0mm}\mbox{\raisebox{-1.3mm}{$\downarrow$}}}
             $\Big\downarrow$}
\def\maponD{\onbelow}
\def\MaponD#1#2{\onbelow\llap{$\vcenter{\hbox{$\Back\!\scriptstyle#1$}}$}
                     \rlap{$\vcenter{\hbox{$\scriptstyle#2$}}$}}

% diagonal maps (SE)
\def\diagSE{\smash{\back\back\back\unitlength=1mm
            \begin{picture}(10,8)\put(0,8){\line(1,-1){10}}\end{picture}}}

\def\mapSE{\smash{\diagSE
                  \rlap{\raisebox{-1.5666mm}{$\rule{-4mm}{0mm}\searrow$}}}}
                                                 
\def\MapSE#1{\smash{\diagSE
             \rlap{\raisebox{-1.5666mm}{$\rule{-4mm}{0mm}\searrow$}}
             \rlap{\raisebox{3mm}{$\rule{-3.5mm}{0mm}
                                  \vcenter{\hbox{$\scriptstyle#1$}}$}}}}

% older maps (delete later?)
\def\mapright{\smash{\mathop{\longrightarrow}}}
\def\mapincl{\smash{\mathop{\hookrightarrow}}}
\def\mapup{\Big\uparrow}
\def\mapdown{\Big\downarrow}
\def\mapdowneq{\Big\parallel}

\def\Mapright#1{\smash{\mathop{\longrightarrow}\limits^{#1}}}
\def\Mapincl#1{\smash{\mathop{\hookrightarrow}\limits^{#1}}}
\def\Mapup#1{\Big\uparrow\rlap{$\vcenter{\hbox{$\scriptstyle#1$}}$}}
\def\Mapdown#1{\Big\downarrow\rlap{$\vcenter{\hbox{$\scriptstyle#1$}}$}}
\def\Mapdowneq#1{\Big\parallel\rlap{$\vcenter{\hbox{$\scriptstyle#1$}}$}}


% lauracode.txt END

% p.tex BEGIN
%% preamble.tex
%% this should be included with a command like
%% %% preamble.tex
%% this should be included with a command like
\usepackage{expl3}
%% %% preamble.tex
%% this should be included with a command like
\usepackage{expl3}
%% %% preamble.tex
%% this should be included with a command like
\usepackage{expl3}
%% \input{./code/p}
\input{./code/mips.sty}
%\usepackage{mips}
\usepackage{amsfonts}
\usepackage{amsthm}
\usepackage{latexsym}
\usepackage{amsmath}
\usepackage{amsmath}				% want AMS fonts
\usepackage{amssymb}                % use AMS symbols
%\usepackage[ruled]{algorithm2e}
\usepackage{enumerate}
\usepackage{enumitem}
\usepackage{setspace}
\usepackage[normalem]{ulem}
% \usepackage{graphicx, subfigure}

\usepackage{multicol}
\usepackage{multirow}
\usepackage{bytefield}
\usepackage{color}
\usepackage{xcolor}
\usepackage{listings}
\usepackage{chronology}
\usepackage{pdfpages}

\usepackage{forloop}
\usepackage{xparse}
\usepackage{xstring}
\usepackage{listofitems}

\usepackage{courier}
\usepackage{textcomp}

\newcommand{\sk}{s}
\newtheorem{theorem}{Theorem}
\newtheorem{lesson}{Lesson}
\newtheorem{proposition}{Proposition}
\newtheorem{lemma}{Lemma}
\newtheorem{corollary}{Corollary}
\newtheorem{fact}{Fact}
\newtheorem*{claim}{Claim}
%\theoremstyle{definition}
\newtheorem{definition}{Definition}
\newtheorem{assumption}{Assumption}
\theoremstyle{remark}
\newtheorem{example}{Example}
\newtheorem*{remark}{Remark}

\DeclareSymbolFont{AMSb}{U}{msb}{m}{n}
\DeclareMathSymbol{\F}{\mathalpha}{AMSb}{"46}
\DeclareMathSymbol{\N}{\mathalpha}{AMSb}{"4E}
\DeclareMathSymbol{\R}{\mathalpha}{AMSb}{"52}
\DeclareMathSymbol{\X}{\mathalpha}{AMSb}{"58}
\DeclareMathSymbol{\Zz}{\mathalpha}{AMSb}{"5A}
\newcommand{\Z}[1]{{\ensuremath{\Zz_{#1}} }}
\newcommand{\Zs}[1]{\ensuremath{\Zz^{\ast}_{#1}}}
\newcommand{\Zn}{\Z{n}}
\newcommand{\Zns}{\Zs{n}}
\newcommand{\Zstar}[1]{\Zs{#1}}
\newcommand{\Zp}{{\Z{p}}}
\newcommand{\Zps}{\Zs{p}}
\newcommand{\Zqs}{\Zs{q}}
\newcommand{\Zq}{{\Z{q}}}
%\newcommand{\QR}{\mathop{\mathrm{QR}}\nolimits}
\newcommand{\ord}[1]{\mathop{\mathrm{ord}}({#1})}
\newcommand{\QR}[1]{\ensuremath{\textit{QR}_{#1}}}
\newcommand{\becomes}{:=}
\newcommand{\rem}[1]{\ensuremath{\ \operatorname{rem} #1}}  
\newcommand{\U}{{\mathcal{U}}}
\newcommand{\floor}[1]{\ensuremath{\lfloor{#1}\rfloor}}
\newcommand{\de}[1]{\ensuremath{\Delta{#1}}}
\newcommand{\js}[2]{\left( \frac{#1}{#2} \right)}

\renewcommand{\QR}{{\mbox{QR}}}
\newcommand{\QNR}{{\mbox{QNR}}}
\newcommand{\crt}{{\mbox{CRT}}}
\newcommand{\rsa}{{\mbox{RSA}}}
\newcommand{\rsamod}{{\mbox{RSA-modulus}}}


\newcommand{\greq}[1]{\stackrel{#1}{=}}
\newcommand{\hash}{\ensuremath{\mathcal{H}}}
\newcommand{\negl}{{\tt neg}}
\newcommand{\cindist}{\stackrel{c}{\approx}}
\newcommand{\A}{{\mathcal{A}}}
\newcommand{\B}{{\mathcal{B}}}

\newcommand{\gen}{\mathit{Gen}}
\newcommand{\keygen}{\mathit{KeyGen}}
\newcommand{\RSA}{\mathit{RSA}}
\newcommand{\pk}{\mathit{pk}}
\newcommand{\PK}{\mathit{PK}}
\newcommand{\SK}{\mathit{SK}}
\newcommand{\sign}{\mathit{Sign}}
\newcommand{\verify}{\mathit{Verify}}


% \setlength{\oddsidemargin}{.25in}
% \setlength{\evensidemargin}{.25in}
% \setlength{\textwidth}{6in}
% \setlength{\topmargin}{-0.4in}
% \setlength{\textheight}{8.5in}

\newcommand{\handout}[5]{
   \renewcommand{\thepage}{#1-\arabic{page}}
   \noindent
   \begin{center}
   \framebox{
      \vbox{
    \hbox to 5.78in { {\bf ECE251: Computer Architecture} \hfill #2 }
       \vspace{4mm}
       \hbox to 5.78in { {\Large \hfill #5  \hfill} }
       \vspace{2mm}
       \hbox to 5.78in { {\it #3 \hfill #4} }
      }
   }
   \end{center}
   \vspace*{3mm}
}

\newcommand{\ho}[4]{\handout{#1}{#2}{Professor:
#3}{}{Handout #1: #4}}

\newcommand{\lnotes}[3]{\handout{#1}{#2}{Professor:
#3}{}{Lecture #1}}

\newcommand{\solution}[1]{#1}
%\homework{number}{out}{due}{instructor}
\newcommand{\homework}[4]{\handout{HW #1}{#2}{Professor: #4}{Due: #3}{Homework Assignment #1}}

%================================================
% problemset macros
%================================================
% count problems
\newcounter{solutioncount}
\setcounter{solutioncount}{0}
\newcommand{\problem}[1]{%
\addtocounter{solutioncount}{1}%
\section*{Problem \arabic{solutioncount}: #1}}

% lets you make alphabetical lists (at first-level of enumeration)
\newenvironment
  {alphabetize}{\renewcommand{\theenumi}{\alph{enumi}}\begin{enumerate}}
  {\end{enumerate}\renewcommand{\theenumi}{\arabic{enumi}}}



\input{./code/mips.sty}
%\usepackage{mips}
\usepackage{amsfonts}
\usepackage{amsthm}
\usepackage{latexsym}
\usepackage{amsmath}
\usepackage{amsmath}				% want AMS fonts
\usepackage{amssymb}                % use AMS symbols
%\usepackage[ruled]{algorithm2e}
\usepackage{enumerate}
\usepackage{enumitem}
\usepackage{setspace}
\usepackage[normalem]{ulem}
% \usepackage{graphicx, subfigure}

\usepackage{multicol}
\usepackage{multirow}
\usepackage{bytefield}
\usepackage{color}
\usepackage{xcolor}
\usepackage{listings}
\usepackage{chronology}
\usepackage{pdfpages}

\usepackage{forloop}
\usepackage{xparse}
\usepackage{xstring}
\usepackage{listofitems}

\usepackage{courier}
\usepackage{textcomp}

\newcommand{\sk}{s}
\newtheorem{theorem}{Theorem}
\newtheorem{lesson}{Lesson}
\newtheorem{proposition}{Proposition}
\newtheorem{lemma}{Lemma}
\newtheorem{corollary}{Corollary}
\newtheorem{fact}{Fact}
\newtheorem*{claim}{Claim}
%\theoremstyle{definition}
\newtheorem{definition}{Definition}
\newtheorem{assumption}{Assumption}
\theoremstyle{remark}
\newtheorem{example}{Example}
\newtheorem*{remark}{Remark}

\DeclareSymbolFont{AMSb}{U}{msb}{m}{n}
\DeclareMathSymbol{\F}{\mathalpha}{AMSb}{"46}
\DeclareMathSymbol{\N}{\mathalpha}{AMSb}{"4E}
\DeclareMathSymbol{\R}{\mathalpha}{AMSb}{"52}
\DeclareMathSymbol{\X}{\mathalpha}{AMSb}{"58}
\DeclareMathSymbol{\Zz}{\mathalpha}{AMSb}{"5A}
\newcommand{\Z}[1]{{\ensuremath{\Zz_{#1}} }}
\newcommand{\Zs}[1]{\ensuremath{\Zz^{\ast}_{#1}}}
\newcommand{\Zn}{\Z{n}}
\newcommand{\Zns}{\Zs{n}}
\newcommand{\Zstar}[1]{\Zs{#1}}
\newcommand{\Zp}{{\Z{p}}}
\newcommand{\Zps}{\Zs{p}}
\newcommand{\Zqs}{\Zs{q}}
\newcommand{\Zq}{{\Z{q}}}
%\newcommand{\QR}{\mathop{\mathrm{QR}}\nolimits}
\newcommand{\ord}[1]{\mathop{\mathrm{ord}}({#1})}
\newcommand{\QR}[1]{\ensuremath{\textit{QR}_{#1}}}
\newcommand{\becomes}{:=}
\newcommand{\rem}[1]{\ensuremath{\ \operatorname{rem} #1}}  
\newcommand{\U}{{\mathcal{U}}}
\newcommand{\floor}[1]{\ensuremath{\lfloor{#1}\rfloor}}
\newcommand{\de}[1]{\ensuremath{\Delta{#1}}}
\newcommand{\js}[2]{\left( \frac{#1}{#2} \right)}

\renewcommand{\QR}{{\mbox{QR}}}
\newcommand{\QNR}{{\mbox{QNR}}}
\newcommand{\crt}{{\mbox{CRT}}}
\newcommand{\rsa}{{\mbox{RSA}}}
\newcommand{\rsamod}{{\mbox{RSA-modulus}}}


\newcommand{\greq}[1]{\stackrel{#1}{=}}
\newcommand{\hash}{\ensuremath{\mathcal{H}}}
\newcommand{\negl}{{\tt neg}}
\newcommand{\cindist}{\stackrel{c}{\approx}}
\newcommand{\A}{{\mathcal{A}}}
\newcommand{\B}{{\mathcal{B}}}

\newcommand{\gen}{\mathit{Gen}}
\newcommand{\keygen}{\mathit{KeyGen}}
\newcommand{\RSA}{\mathit{RSA}}
\newcommand{\pk}{\mathit{pk}}
\newcommand{\PK}{\mathit{PK}}
\newcommand{\SK}{\mathit{SK}}
\newcommand{\sign}{\mathit{Sign}}
\newcommand{\verify}{\mathit{Verify}}


% \setlength{\oddsidemargin}{.25in}
% \setlength{\evensidemargin}{.25in}
% \setlength{\textwidth}{6in}
% \setlength{\topmargin}{-0.4in}
% \setlength{\textheight}{8.5in}

\newcommand{\handout}[5]{
   \renewcommand{\thepage}{#1-\arabic{page}}
   \noindent
   \begin{center}
   \framebox{
      \vbox{
    \hbox to 5.78in { {\bf ECE251: Computer Architecture} \hfill #2 }
       \vspace{4mm}
       \hbox to 5.78in { {\Large \hfill #5  \hfill} }
       \vspace{2mm}
       \hbox to 5.78in { {\it #3 \hfill #4} }
      }
   }
   \end{center}
   \vspace*{3mm}
}

\newcommand{\ho}[4]{\handout{#1}{#2}{Professor:
#3}{}{Handout #1: #4}}

\newcommand{\lnotes}[3]{\handout{#1}{#2}{Professor:
#3}{}{Lecture #1}}

\newcommand{\solution}[1]{#1}
%\homework{number}{out}{due}{instructor}
\newcommand{\homework}[4]{\handout{HW #1}{#2}{Professor: #4}{Due: #3}{Homework Assignment #1}}

%================================================
% problemset macros
%================================================
% count problems
\newcounter{solutioncount}
\setcounter{solutioncount}{0}
\newcommand{\problem}[1]{%
\addtocounter{solutioncount}{1}%
\section*{Problem \arabic{solutioncount}: #1}}

% lets you make alphabetical lists (at first-level of enumeration)
\newenvironment
  {alphabetize}{\renewcommand{\theenumi}{\alph{enumi}}\begin{enumerate}}
  {\end{enumerate}\renewcommand{\theenumi}{\arabic{enumi}}}



\input{./code/mips.sty}
%\usepackage{mips}
\usepackage{amsfonts}
\usepackage{amsthm}
\usepackage{latexsym}
\usepackage{amsmath}
\usepackage{amsmath}				% want AMS fonts
\usepackage{amssymb}                % use AMS symbols
%\usepackage[ruled]{algorithm2e}
\usepackage{enumerate}
\usepackage{enumitem}
\usepackage{setspace}
\usepackage[normalem]{ulem}
% \usepackage{graphicx, subfigure}

\usepackage{multicol}
\usepackage{multirow}
\usepackage{bytefield}
\usepackage{color}
\usepackage{xcolor}
\usepackage{listings}
\usepackage{chronology}
\usepackage{pdfpages}

\usepackage{forloop}
\usepackage{xparse}
\usepackage{xstring}
\usepackage{listofitems}

\usepackage{courier}
\usepackage{textcomp}

\newcommand{\sk}{s}
\newtheorem{theorem}{Theorem}
\newtheorem{lesson}{Lesson}
\newtheorem{proposition}{Proposition}
\newtheorem{lemma}{Lemma}
\newtheorem{corollary}{Corollary}
\newtheorem{fact}{Fact}
\newtheorem*{claim}{Claim}
%\theoremstyle{definition}
\newtheorem{definition}{Definition}
\newtheorem{assumption}{Assumption}
\theoremstyle{remark}
\newtheorem{example}{Example}
\newtheorem*{remark}{Remark}

\DeclareSymbolFont{AMSb}{U}{msb}{m}{n}
\DeclareMathSymbol{\F}{\mathalpha}{AMSb}{"46}
\DeclareMathSymbol{\N}{\mathalpha}{AMSb}{"4E}
\DeclareMathSymbol{\R}{\mathalpha}{AMSb}{"52}
\DeclareMathSymbol{\X}{\mathalpha}{AMSb}{"58}
\DeclareMathSymbol{\Zz}{\mathalpha}{AMSb}{"5A}
\newcommand{\Z}[1]{{\ensuremath{\Zz_{#1}} }}
\newcommand{\Zs}[1]{\ensuremath{\Zz^{\ast}_{#1}}}
\newcommand{\Zn}{\Z{n}}
\newcommand{\Zns}{\Zs{n}}
\newcommand{\Zstar}[1]{\Zs{#1}}
\newcommand{\Zp}{{\Z{p}}}
\newcommand{\Zps}{\Zs{p}}
\newcommand{\Zqs}{\Zs{q}}
\newcommand{\Zq}{{\Z{q}}}
%\newcommand{\QR}{\mathop{\mathrm{QR}}\nolimits}
\newcommand{\ord}[1]{\mathop{\mathrm{ord}}({#1})}
\newcommand{\QR}[1]{\ensuremath{\textit{QR}_{#1}}}
\newcommand{\becomes}{:=}
\newcommand{\rem}[1]{\ensuremath{\ \operatorname{rem} #1}}  
\newcommand{\U}{{\mathcal{U}}}
\newcommand{\floor}[1]{\ensuremath{\lfloor{#1}\rfloor}}
\newcommand{\de}[1]{\ensuremath{\Delta{#1}}}
\newcommand{\js}[2]{\left( \frac{#1}{#2} \right)}

\renewcommand{\QR}{{\mbox{QR}}}
\newcommand{\QNR}{{\mbox{QNR}}}
\newcommand{\crt}{{\mbox{CRT}}}
\newcommand{\rsa}{{\mbox{RSA}}}
\newcommand{\rsamod}{{\mbox{RSA-modulus}}}


\newcommand{\greq}[1]{\stackrel{#1}{=}}
\newcommand{\hash}{\ensuremath{\mathcal{H}}}
\newcommand{\negl}{{\tt neg}}
\newcommand{\cindist}{\stackrel{c}{\approx}}
\newcommand{\A}{{\mathcal{A}}}
\newcommand{\B}{{\mathcal{B}}}

\newcommand{\gen}{\mathit{Gen}}
\newcommand{\keygen}{\mathit{KeyGen}}
\newcommand{\RSA}{\mathit{RSA}}
\newcommand{\pk}{\mathit{pk}}
\newcommand{\PK}{\mathit{PK}}
\newcommand{\SK}{\mathit{SK}}
\newcommand{\sign}{\mathit{Sign}}
\newcommand{\verify}{\mathit{Verify}}


% \setlength{\oddsidemargin}{.25in}
% \setlength{\evensidemargin}{.25in}
% \setlength{\textwidth}{6in}
% \setlength{\topmargin}{-0.4in}
% \setlength{\textheight}{8.5in}

\newcommand{\handout}[5]{
   \renewcommand{\thepage}{#1-\arabic{page}}
   \noindent
   \begin{center}
   \framebox{
      \vbox{
    \hbox to 5.78in { {\bf ECE251: Computer Architecture} \hfill #2 }
       \vspace{4mm}
       \hbox to 5.78in { {\Large \hfill #5  \hfill} }
       \vspace{2mm}
       \hbox to 5.78in { {\it #3 \hfill #4} }
      }
   }
   \end{center}
   \vspace*{3mm}
}

\newcommand{\ho}[4]{\handout{#1}{#2}{Professor:
#3}{}{Handout #1: #4}}

\newcommand{\lnotes}[3]{\handout{#1}{#2}{Professor:
#3}{}{Lecture #1}}

\newcommand{\solution}[1]{#1}
%\homework{number}{out}{due}{instructor}
\newcommand{\homework}[4]{\handout{HW #1}{#2}{Professor: #4}{Due: #3}{Homework Assignment #1}}

%================================================
% problemset macros
%================================================
% count problems
\newcounter{solutioncount}
\setcounter{solutioncount}{0}
\newcommand{\problem}[1]{%
\addtocounter{solutioncount}{1}%
\section*{Problem \arabic{solutioncount}: #1}}

% lets you make alphabetical lists (at first-level of enumeration)
\newenvironment
  {alphabetize}{\renewcommand{\theenumi}{\alph{enumi}}\begin{enumerate}}
  {\end{enumerate}\renewcommand{\theenumi}{\arabic{enumi}}}



\input{./code/mips.sty}
%\usepackage{mips}
\usepackage{amsfonts}
\usepackage{amsthm}
\usepackage{latexsym}
\usepackage{amsmath}
\usepackage{expl3}
\usepackage{amsmath}				% want AMS fonts
\usepackage{amssymb}                            % use AMS symbols
%\usepackage[ruled]{algorithm2e}
\usepackage{enumerate}
\usepackage{setspace}
\usepackage[normalem]{ulem}
% \usepackage{graphicx, subfigure}

\usepackage{multicol}
\usepackage{multirow}
\usepackage{bytefield}
\usepackage{color}
\usepackage{xcolor}
\usepackage{listings}
\usepackage{chronology}
\usepackage{pdfpages}

\usepackage{forloop}
\usepackage{xparse}
\usepackage{xstring}
\usepackage{listofitems}

\usepackage{courier}
\usepackage{textcomp}
\newcommand{\colorbitbox}[4][lrbt]{%
    \rlap{\bitbox[#1]{#3}{\color{#2}\rule{\dimexpr\width-0.4pt}{\dimexpr\height-0.4pt}}}%
    \bitbox[#1]{#3}{#4}}

\definecolor{lightgray}{gray}{0.85}

\newcommand{\sk}{s}
\newtheorem{theorem}{Theorem}
\newtheorem{lesson}{Lesson}
\newtheorem{proposition}{Proposition}
\newtheorem{lemma}{Lemma}
\newtheorem{corollary}{Corollary}
\newtheorem{fact}{Fact}
\newtheorem*{claim}{Claim}
%\theoremstyle{definition}
\newtheorem{definition}{Definition}
\newtheorem{assumption}{Assumption}
\theoremstyle{remark}
\newtheorem{example}{Example}
\newtheorem*{remark}{Remark}

\DeclareSymbolFont{AMSb}{U}{msb}{m}{n}
\DeclareMathSymbol{\F}{\mathalpha}{AMSb}{"46}
\DeclareMathSymbol{\N}{\mathalpha}{AMSb}{"4E}
\DeclareMathSymbol{\R}{\mathalpha}{AMSb}{"52}
\DeclareMathSymbol{\X}{\mathalpha}{AMSb}{"58}
\DeclareMathSymbol{\Zz}{\mathalpha}{AMSb}{"5A}
\newcommand{\Z}[1]{{\ensuremath{\Zz_{#1}} }}
\newcommand{\Zs}[1]{\ensuremath{\Zz^{\ast}_{#1}}}
\newcommand{\Zn}{\Z{n}}
\newcommand{\Zns}{\Zs{n}}
\newcommand{\Zstar}[1]{\Zs{#1}}
\newcommand{\Zp}{{\Z{p}}}
\newcommand{\Zps}{\Zs{p}}
\newcommand{\Zqs}{\Zs{q}}
\newcommand{\Zq}{{\Z{q}}}
%\newcommand{\QR}{\mathop{\mathrm{QR}}\nolimits}
\newcommand{\ord}[1]{\mathop{\mathrm{ord}}({#1})}
\newcommand{\QR}[1]{\ensuremath{\textit{QR}_{#1}}}
\newcommand{\becomes}{:=}
\newcommand{\rem}[1]{\ensuremath{\ \operatorname{rem} #1}}  
\newcommand{\U}{{\mathcal{U}}}
\newcommand{\floor}[1]{\ensuremath{\lfloor{#1}\rfloor}}
\newcommand{\de}[1]{\ensuremath{\Delta{#1}}}
\newcommand{\js}[2]{\left( \frac{#1}{#2} \right)}

\renewcommand{\QR}{{\mbox{QR}}}
\newcommand{\QNR}{{\mbox{QNR}}}
\newcommand{\crt}{{\mbox{CRT}}}
\newcommand{\rsa}{{\mbox{RSA}}}
\newcommand{\rsamod}{{\mbox{RSA-modulus}}}


\newcommand{\greq}[1]{\stackrel{#1}{=}}
\newcommand{\hash}{\ensuremath{\mathcal{H}}}
\newcommand{\negl}{{\tt neg}}
\newcommand{\cindist}{\stackrel{c}{\approx}}
\newcommand{\A}{{\mathcal{A}}}
\newcommand{\B}{{\mathcal{B}}}

\newcommand{\gen}{\mathit{Gen}}
\newcommand{\keygen}{\mathit{KeyGen}}
\newcommand{\RSA}{\mathit{RSA}}
\newcommand{\pk}{\mathit{pk}}
\newcommand{\PK}{\mathit{PK}}
\newcommand{\SK}{\mathit{SK}}
\newcommand{\sign}{\mathit{Sign}}
\newcommand{\verify}{\mathit{Verify}}


\setlength{\oddsidemargin}{.25in}
\setlength{\evensidemargin}{.25in}
\setlength{\textwidth}{6in}
\setlength{\topmargin}{-0.4in}
\setlength{\textheight}{8.5in}

\newcommand{\handout}[5]{
   \renewcommand{\thepage}{#1-\arabic{page}}
   \noindent
   \begin{center}
   \framebox{
      \vbox{
    \hbox to 5.78in { {\bf ECE251: Computer Architecture} \hfill #2 }
       \vspace{4mm}
       \hbox to 5.78in { {\Large \hfill #5  \hfill} }
       \vspace{2mm}
       \hbox to 5.78in { {\it #3 \hfill #4} }
      }
   }
   \end{center}
   \vspace*{3mm}
}

\newcommand{\ho}[4]{\handout{#1}{#2}{Professor:
#3}{}{Handout #1: #4}}

\newcommand{\lnotes}[3]{\handout{#1}{#2}{Professor:
#3}{}{Lecture #1}}

\newcommand{\solution}[1]{#1}
%\homework{number}{out}{due}{instructor}
\newcommand{\homework}[4]{\handout{HW #1}{#2}{Professor: #4}{Due: #3}{Homework Assignment #1}}

%================================================
% problemset macros
%================================================
% count problems
\newcounter{solutioncount}
\setcounter{solutioncount}{0}
\newcommand{\problem}[1]{%
\addtocounter{solutioncount}{1}%
\section*{Problem \arabic{solutioncount}: #1}}

% lets you make alphabetical lists (at first-level of enumeration)
\newenvironment
  {alphabetize}{\renewcommand{\theenumi}{\alph{enumi}}\begin{enumerate}}
  {\end{enumerate}\renewcommand{\theenumi}{\arabic{enumi}}}


%
% \ExplSyntaxOn
% \NewDocumentCommand{\loopBin}{m}
%  {
%   % save the input in a variable
%   \tl_set:Nn \mips_getletters { #1 }
%   % replace spaces with \textvisiblespace
%   \tl_replace_all:Nnn \mips_getletters { ~ } { \textvisiblespace }
%   % map the input surrounding each item by parentheses
%   \tl_map_inline:Nn \mips_getletters { ##1 }
%  }
% \tl_new:N \mips_getletters
% \ExplSyntaxOff
%


%
% \newcommand{\binLoop}[1]
% {
%     \begin{bytefield}[boxformatting={\centering\itshape},bitwidth=1.5em,endianness=big]{32}
%     \bitheader{0-31}
%     % \StrLen{#1}
%     \foreach \x in {1,...\StrLen{#1}}
%     {
%         \bitbox{\x,#1{\x}}
%     }
    
%     \end{bytefield}
% }
%

%
 % \NewDocumentCommand{\extract}{mm}
\ExplSyntaxOn
\cs_new:Npn \extract #1 #2
 {
  \str_item:fn { #1 } { #2+1 } % since it starts counting from 1
 }
\cs_generate_variant:Nn \str_item:nn {f}
\ExplSyntaxOff

\ExplSyntaxOn
% \cs_new:Npn \makebitboxes #1 #2
 \NewDocumentCommand{\makebitboxes}{mm}
 {
    \begin{bytefield}[boxformatting={\centering\itshape},bitwidth=1.5em,endianness=big]{\numbits}
    \bitheader{0-{#1 - 1}} \\
    
    \int_step_variable:nnnNn {0}{1}{#1-1}{\tvar}
   {
        % (\tvar,\extract{#2}{\tvar})
        \int_compare:nNnTF { \int_mod:nn {\tvar} { 2 } } = { 0 }
       {% odd box
            \colorbitbox{lightgray}{1}{\extract{#2}{\tvar}}
       }
       { %even box
            \bitbox{1}{\extract{#2}{\tvar}}
       }
   }
   \end{bytefield}
  }
\ExplSyntaxOff
%

\ExplSyntaxOn
% \cs_new:Npn \makebitboxes #1 #2
 \NewDocumentCommand{\makeemptybyteboxes}{m}
 {
    \begin{bytefield}[boxformatting={\centering\itshape},bitwidth=1.5em,endianness=big]{32}
    
    \int_step_variable:nnnNn {0}{1}{#1-1}{\tvar}
   {
        % (\tvar,\extract{#2}{\tvar})
        \int_compare:nNnTF { \int_mod:nn {\tvar} { 2 } } = { 0 }
       {% odd box
            \colorbitbox{lightgray}{4}{}
       }
       { %even box
            \bitbox{4}{}
       }
   }
   \end{bytefield}
  }
\ExplSyntaxOff
%

\ExplSyntaxOn
 \NewDocumentCommand{\makeemptybitboxes}{m}
 {
    \begin{bytefield}[boxformatting={\centering\itshape},bitwidth=1.5em,endianness=big]{\numbits}
    \bitheader{0-{#1 - 1}} \\
    
    \int_step_variable:nnnNn {0}{1}{#1-1}{\tvar}
   {
        % (\tvar,\extract{#2}{\tvar})
        \int_compare:nNnTF { \int_mod:nn {\tvar} { 2 } } = { 0 }
       {% odd box
            \colorbitbox{lightgray}{1}{}
       }
       { %even box
            \bitbox{1}{}
       }
   }
   \end{bytefield}
  }
\ExplSyntaxOff
%
% p.tex END

% testpoints.tex BEGIN
%------------------------------------------------------------------
% PROBLEM, PART, AND POINT COUNTING...

% Create the problem number counter.  Initialize to zero.
\newcounter{problemnum}

% Specify that problems should be labeled with arabic numerals.
\renewcommand{\theproblemnum}{\arabic{problemnum}}


% Create the part-within-a-problem counter, "within" the problem counter.
% This counter resets to zero automatically every time the PROBLEMNUM counter
% is incremented.
\newcounter{partnum}[problemnum]

% Specify that parts should be labeled with lowercase letters.
\renewcommand{\thepartnum}{\alph{partnum}}

% Make a counter to keep track of total points assigned to problems...
\newcounter{totalpoints}

% Make counters to keep track of points for parts...
\newcounter{curprobpts}		% Points assigned for the problem as a whole.
\newcounter{totalparts}		% Total points assigned to the various parts.

% Make a counter to keep track of the number of points on each page...
\newcounter{pagepoints}
% This counter is reset each time a page is printed.

% This "program" keeps track of how many points appear on each page, so that
% the total can be printed on the page itself.  Points are added to the total
% for a page when the PART (not the problem) they are assigned to is specified.
% When a problem without parts appears, the PAGEPOINTS are incremented directly
% from the problem as a whole (CURPROBPTS).


%---------------------------------------------------------------------------


% The \problem environment first checks the information about the previous
% problem.  If no parts appeared (or if they were all assigned zero points,
% then it increments TOTALPOINTS directly from CURPROBPTS, the points assigned
% to the last problem as a whole.  If the last problem did contain parts, it
% checks to make sure that their point values total up to the correct sum.
% It then puts the problem number on the page, along with the points assigned
% to it.

\newenvironment{problem}[1]{
% STATEMENTS TO BE EXECUTED WHEN A NEW PROBLEM IS BEGUN:
%
% Increment the problem number counter, and set the current \ref value to that
% number.
\refstepcounter{problemnum}
%
% Add some vspace to separate from the last problem.
\vspace{0.15in} \par
%
\setcounter{curprobpts}{#1} \setcounter{totalparts}{0}	% Reset counters.
%
% Now put in the "announcement" on the page.
{\Large \bf \theproblemnum. \normalsize ({\it \arabic{curprobpts} point\null\ifnum \value{curprobpts} = 1\else s\fi}\/)}
}{
% STATEMENTS TO BE EXECUTED WHEN AN OLD PROBLEM IS ENDED:
%
% If no parts to problem, then increment TOTALPOINTS and PAGEPOINTS for the
% entire problem at once.
\ifnum \value{totalparts} = 0
	\addtocounter{totalpoints}{\value{curprobpts}}	% Add pts to total.
	\addtocounter{pagepoints}{\value{curprobpts}}	% Add pts to page total.
%
% If there were parts for the problem, then check to make sure they total up
% to the same number of points that the problem is worth. Issue a warning
% if not.
\else \ifnum \value{totalparts} = \value{curprobpts}
	\else \typeout{}
	\typeout{!!!!!!!   POINT ACCOUNTING ERROR   !!!!!!!!}
	\typeout{PROBLEM [\theproblemnum] WAS ALLOCATED \arabic{curprobpts} POINTS,}
	\typeout{BUT CONTAINS PARTS TOTALLING \arabic{totalparts} POINTS!}
	\typeout{}
	\fi
\fi
}


%---------------------------------------------------------------------------


% The \newpart command increments the part counter and displays an appropriate
% lowercase letter to mark the part.  It adds points to the point counter
% immediately.  If 0 points are specified, no point announcement is made.
% Otherwise, the announcement is in scriptsize italics.

\newcommand{\newpart}[1]
{
\refstepcounter{partnum}	% Set the current \ref value to the part number.
\hspace{0.25in}		% Indent the part by a quarter inch.
%
% If points are to be printed for this problem (signaled by point value > 0),
% then put them in in scriptsize italics.
\ifnum #1 > 0
	\makebox[0.5in][l]{{\bf \thepartnum.} {\bf ({\it #1 pt\ifnum #1 = 1\else s\fi\/}) \,\,}}
\else
	\makebox[0.25in][l]{({\bf \thepartnum})}
\fi
%
\hspace{0.1in}		% Lead the material away from the part "number".
%
\addtocounter{totalparts}{#1}	% Add points to totalparts for this problem.
\addtocounter{pagepoints}{#1}	% Add points to total for this page.
\addtocounter{totalpoints}{#1}	% Add points to total for entire test.
}


%---------------------------------------------------------------------------



% Just in case you want to skip some numbers in your test...

\newcommand{\skipproblem}[1]{\addtocounter{problemnum}{#1}}



%---------------------------------------------------------------------------


% The \showpoints command simply gives a count of the total points read in up to
% the location at which the command is placed.  Typically, one places one
% \showpoints command at the end of the latex file, just prior to the
% \end{document} command.  It can appear elsewhere, however.

\newcommand{\showpoints}
{
\typeout{}  
\typeout{====> A TOTAL OF \arabic{totalpoints} POINTS WERE READ.}
\typeout{}
}


%---------------------------------------------------------------------------


% testpoints.tex END